\section{ماژول فشرده‌سازی برش کانال‌های ترانسفورمر (\lr{SliceGPT})}
\label{sec:slicegpt}

\subsection{معرفی و انگیزه}
این مکانیزم بر مبنای مبانی نظری پیشرفته در ادبیات توجه و فشرده‌سازی \cite{ashkboos2024slicegpt} توسعه یافته است.

یکی از نوآورانه‌ترین رویکردها در فشرده‌سازی ساختاریافته مدل‌های زبانی بزرگ، حذف فیزیکی سطرها و ستون‌های ماتریس‌های وزن بدون نیاز به ساختارهای تنک یا کرنل‌های محاسباتی خاص است. روش \lr{SliceGPT} بر مبنای یک ویژگی ریاضی بنیادین در معماری‌های ترانسفورمر پایه‌گذاری شده است: «نامتغیری محاسبات نسبت به تبدیل یکامتعامد در فضای مختصات پنهان» (\lr{Orthogonal Coordinate Invariance}).

در شبکه‌های استاندارد، فعال‌سازی‌ها معمولاً در امتداد جهت‌های متعامد اصلی متمرکز نبوده و در کل فضای مختصات پخش شده‌اند؛ در نتیجه حذف مستقیم سطرها یا ستون‌ها موجب تخریب شدید دقت می‌گردد. الگوریتم \lr{SliceGPT} با چرخش فضای مختصات فعال‌سازی‌ها با یک ماتریس یکامتعامد محاسبه‌شده از کوواریانس فعال‌سازی‌ها، انرژی سیگنال را بر روی تعداد کمتری از مؤلفه‌های اصلی متمرکز می‌کند. پس از این چرخش، مؤلفه‌های با واریانس بسیار اندک را می‌توان به طور فیزیکی حذف (\lr{Slice}) کرد، بدون آنکه خلل قابل توجهی در بازنمایی‌های معنایی مدل پدید آید.

\subsection{مبانی نظری و فرمول‌بندی ریاضی}
اصول ریاضی \lr{SliceGPT} در سه گام متوالی تدوین می‌گردد:

\begin{enumerate}
    \item \textbf{اصل نامتغیری مختصات یکامتعامد:}
    اگر $\mathbf{Q} \in \mathbb{R}^{d_{\text{in}} \times d_{\text{in}}}$ یک ماتریس یکامتعامد باشد به نحوی که $\mathbf{Q} \mathbf{Q}^\top = \mathbf{I}$، ضرب ماتریس وزن $\mathbf{W} \in \mathbb{R}^{d_{\text{out}} \times d_{\text{in}}}$ در فعال‌سازی ورودی $\mathbf{X}$ را می‌توان بدون تغییر در مقدار خروجی بازنویسی نمود:
    \begin{equation}
    \mathbf{W} \mathbf{X} = (\mathbf{W} \mathbf{Q}) (\mathbf{Q}^\top \mathbf{X}) = \tilde{\mathbf{W}} \tilde{\mathbf{X}}
    \end{equation}
    که در آن $\tilde{\mathbf{W}} = \mathbf{W} \mathbf{Q}$ ماتریس وزن چرخانده‌شده و $\tilde{\mathbf{X}} = \mathbf{Q}^\top \mathbf{X}$ فعال‌سازی‌های نگاشت‌شده به دستگاه مختصات جدید هستند.

    \item \textbf{محاسبه ماتریس کوواریانس و تجزیه ویژه‌مقدار (\lr{PCA Rotation}):}
    با استفاده از تانسورهای کالیبراسیون، ماتریس کوواریانس تجربی فعال‌سازی‌ها استخراج می‌شود:
    \begin{equation}
    \mathbf{C} = \frac{1}{N} \mathbf{X} \mathbf{X}^\top = \mathbf{Q} \mathbf{\Lambda} \mathbf{Q}^\top
    \end{equation}
    که در آن $\mathbf{Q}$ ماتریس بردارهای ویژه و $\mathbf{\Lambda}$ ماتریس قطری ویژه‌مقادیر (واریانس در امتداد هر محور) است. ستون‌های $\mathbf{Q}$ به ترتیب نزولی واریانس مرتب می‌گردند.

    \item \textbf{برش ابعادی و پروجکشن فیزیکی (\lr{Subspace Slicing}):}
    با انتخاب ابعاد برش‌یافته $r = \lfloor d_{\text{in}} \cdot (1 - \alpha) \rfloor$ که در آن $\alpha$ نسبت برش کانال‌ها است، ماتریس پروجکشن برش $\mathbf{P}_r \in \mathbb{R}^{d_{\text{in}} \times r}$ با انتخاب $r$ ستون نخست ماتریس یکه تشکیل می‌شود:
    \begin{equation}
    \mathbf{P}_r = [\mathbf{I}_r, \mathbf{0}]^\top
    \end{equation}
    ماتریس وزن جدید به صورت فیزیکی به ابعاد کوچکتری فشرده می‌گردد:
    \begin{equation}
    \mathbf{W}_{\text{sliced}} = \tilde{\mathbf{W}} \mathbf{P}_r = \mathbf{W} \mathbf{Q} \mathbf{P}_r \in \mathbb{R}^{d_{\text{out}} \times r}
    \end{equation}
    این تبدیل اندازه بعد پنهان شبکه را در عمل به $r$ کاهش می‌دهد.
\end{enumerate}

\subsection{معماری و پیاده‌سازی}
پیاده‌سازی الگوریتم \lr{SliceGPT} در مسیر \lr{neuralops/compression\_engine/slicegpt.py} پیاده‌سازی شده و تحت شناسه \lr{slicegpt} در رجیستری فریم‌ورک ثبت گردیده است. مؤلفه‌های کلیدی این پیاده‌سازی عبارتند از:
\begin{itemize}
    \item \textbf{جمع‌آوری هوشمند فعال‌سازی‌ها با قلاب‌های محاسباتی (\lr{Forward Hooks}):} ماژول با تزریق قلاب به لایه‌های هدف، فعال‌سازی‌های واقعی تولیدشده بر روی داده‌های کالیبراسیون را ذخیره و ماتریس کوواریانس را استخراج می‌نماید.
    \item \textbf{سازگاری کامل با لایه‌های \lr{Linear} و \lr{Conv1D}:} انطباق کامل با ساختارهای کانوولوشن ۱ بعدی از طریق وارسی ویژگی‌های کلاس انجام می‌پذیرد.
    \item \textbf{سنجش خطای بازسازی نسبی و تنزیل اضطراری:} پس از چرخش و برش، خطای فروبنیوس ماتریس خروجی نسبت به فعال‌سازی‌های مبنا ارزیابی می‌شود؛ در صورت تخطی از حد مجاز (\lr{max\_error})، پارامتر \lr{fallback\_to\_original} فعال شده و لایه دست‌نخورده باقی می‌ماند.
    \item \textbf{پشتیبانی از لایه‌های نرمال‌سازی (\lr{LayerNorm}):} اعمال ماتریس چرخش $\mathbf{Q}$ به پارامترهای تبدیل لایه‌های نرمال‌سازی برای حفظ پیوستگی عددی میان بلوک‌های پیاپی ترانسفورمر.
\end{itemize}

\subsection{نتایج تجربی}
عملکرد ماژول \lr{SliceGPT} در آزمون‌های بنچ‌مارک ارزیابی چندمدلی در جدول \ref{tab:benchmark_slicegpt} خلاصه شده است:

\begin{table}[H]
\centering
\small
\caption{نتایج بنچ‌مارک تجربی ماژول فشرده‌ساز برش ابعادی \lr{SliceGPT}}
\label{tab:benchmark_slicegpt}
\begin{tabular}{lccccc}
\toprule
\textbf{مدل پایه \lr{(Base Model)}} & \textbf{دقت \lr{(Accuracy)}} & \textbf{شاخص \lr{F1}} & \textbf{حافظه \lr{(Peak VRAM)}} & \textbf{توان \lr{(Throughput)}} & \textbf{وضعیت} \\
\midrule
\lr{DistilBERT} & \lr{80.62\%} & \lr{0.8312} & \lr{339.71 MB} & \lr{126,869.88 tok/s} & \lr{PASS} \\
\lr{RoBERTa} & \lr{50.92\%} & \lr{0.6748} & \lr{366.79 MB} & \lr{84,266.24 tok/s} & \lr{PASS\_LOW} \\
\lr{GPT-2} & \lr{59.06\%} & \lr{0.6270} & \lr{637.16 MB} & \lr{62,933.00 tok/s} & \lr{PASS} \\
\lr{Qwen} & -- & -- & -- & -- & \lr{SKIP} \\
\bottomrule
\end{tabular}
\end{table}

تحلیل شواهد تجربی ثبت‌شده نشان می‌دهد:
\begin{itemize}
    \item در مدل \lr{DistilBERT}، اعمال برش به دقت رضایت‌بخش \lr{80.62\%} با شاخص \lr{F1} معادل \lr{0.8312} و توان پردازشی بسیار بالای \lr{126,869.88 tok/s} انجامیده و معیار قبولی \lr{PASS} حاصل گردیده است.
    \item در مدل \lr{GPT-2}، دقت \lr{59.06\%} با افت ناچیز \lr{1.72\%} نسبت به مدل پایه حفظ گردیده و توان محاسباتی به میزان \lr{62,933.00 tok/s} به همراه وضعیت \lr{PASS} ثبت شده است.
    \item در مدل \lr{RoBERTa}، به دلیل حساسیت اتصالات پسماند به چرخش‌های مختصات و عدم بازتنظیم پس از برش، دقت مدل به \lr{50.92\%} تنزل یافته و وضعیت \lr{PASS\_LOW} ثبت گردیده است.
    \item اجرای مدل \lr{Qwen} به دلیل الزامات منابع سخت‌افزاری در این فاز در وضعیت \lr{SKIP} قرار گرفته است.
\end{itemize}

\subsection{محدودیت‌ها و تحلیل شکست}
محدودیت‌های فنی و عملیاتی ماژول \lr{SliceGPT} به شرح زیر است:
\begin{itemize}
    \item \textbf{وابستگی شدید به کیفیت داده‌های کالیبراسیون:} چرخش بهینه مختصات مستقیماً وابسته به ماتریس کوواریانس فعال‌سازی‌ها است. در صورتی که داده‌های کالیبراسیون تنوع زبانی کافی نداشته باشند، جهات با واریانس بالا به اشتباه شناسایی شده و برش ابعادی موجب افت عملکرد مدل در دامنه‌های جدید می‌گردد.
    \item \textbf{ناهماهنگی در شبکه‌های با اتصالات پسماند متقاطع (\lr{Residual Connections}):} تغییر بعد پنهان در یک بلوک ترانسفورمر مستلزم تطبیق بعد در تمامی اتصالات پسماند موازی است؛ هرگونه ناهمگونی در برش کانال‌ها میان لایه‌های متوالی می‌تواند موجب خطای انطباق ابعاد تانسور یا افت شدید عملکرد (\lr{PASS\_LOW}) شود.
\end{itemize}
